\section{Related Literature}\label{sec:literature}

\paragraph{The joint labor-market risk of spouses.} That spouses' outcomes are linked is known. Their earnings covary,
mostly through permanent components \citep{hyslop2001rising}, more so in booms \citep{shore2010better}; couples sort on
employer and industry, and about a tenth of dual-earner couples work at the same establishment
\citep{hyatt2019coworking}; displacement studies observe couples displaced by the same plant event and drop them
\citep{halla2020job}; and correlated unemployment has been offered as one of two explanations for the small
added-worker response in the United States \citep{kawano2025who}. What has not been measured is the joint job loss
itself: how much a husband's job loss raises his wife's own job-loss hazard, when, and through which channel.
Table~\ref{tab:contribution} sets out what is known and what this paper adds.

\paragraph{Insurance within the household.} Wives' earnings fall by up to 73 cents for each dollar of unemployment
insurance their husbands receive \citep{cullen2000does}. In the PSID, wives' labor supply rises after a husband's
displacement, slowly, replacing more than a quarter of his lost income in the long run \citep{stephens2002worker};
family labor supply is a main channel of insurance against permanent wage shocks \citep{blundell2016consumption}. In
Austria, wives replace a very small share of the household's loss, public transfers and taxes do more, and young
children shape the response \citep{halla2020job}. \citet{birinci2019spousal} attributes the small US response to
crowding out by transfers and, in a model with a frictional labor market and family labor supply, finds that optimal
policy shifts transfers from means-tested to employment-tested; \citet{mankart2017household} also model household
search as insurance. \citet{bacher2025joint} show in CPS couples that a non-participating spouse is 5.6 percentage points more likely to
enter the labor force when the primary earner loses his job, 7.1 points at ages 26--35 and 1.4 at 56--65, and explain
the decline with age in a monthly life-cycle model of joint search with assets, human capital and unemployment
insurance; separation shocks in their model are independent across spouses, and they note that correlated shocks,
from sorting into the same industries or occupations, are left out. This literature treats the wife's job as given
when the husband's is lost. This paper adds the
second earner's own risk and its timing.

\paragraph{Why another paper after \citet{blundell2016consumption}.} The closest paper is
\citet{blundell2016consumption}, who estimate a life-cycle model of two earners with correlated permanent and
transitory wage shocks, nonseparable preferences and a progressive tax and transfer schedule that includes the Earned
Income Tax Credit and food stamps, on PSID consumption, earnings and assets for 1999--2009. A 10\% permanent fall in the
husband's wage lowers consumption by 3.2\%; of the 39 percentage points of consumption that are insured, 63\% are insured
by family labor supply, 20\% by taxes and transfers and 17\% by savings, and a parameter for insurance outside these
channels is not different from zero ($-0.137$, s.e.\ 0.270). The spouses' transitory wage shocks are correlated (0.31),
their permanent shocks are not (covariance 0.0027, s.e.\ 0.0023). Read at face value, the wife is the household's main
insurer; this paper finds that, against job loss, she largely is not.

I rerun their code on their PSID panel (Appendix~\ref{app:bps}). Every point estimate is reproduced
(Table~\ref{tab:bps_relax}, column 1), and the result does not rest on most of their sample choices. Their
controls for both spouses' employment status do not affect the accounting, although they matter at the event itself
(the wife's earnings residual in the interval of her own displacement is $-0.10$ with them and $-0.20^{***}$ without);
adding the PSID's low-income sample raises the share of family labor supply to 0.68. The authors note that dropping
husbands with long spells without work may understate or overstate added-worker effects
\citep[p.~403]{blundell2016consumption}; dropping only the year without work, rather than the husband's whole record,
measures this and lowers the family share from 0.62 to 0.56. The estimator itself is sensitive to the trim of years of
collapse and recovery, a standard treatment of measurement error, which removes almost no displacement years.\footnote{Code:
\texttt{IGE/bps\_reconciliation/} (copies of their do-files with every edit listed in \texttt{README.md}). The 1999
early-release income file needed to rebuild their panel from the raw PSID is not public, so the replication starts from
their constructed panel, which their own code reproduces exactly.}

The difference lies in the shock and in how often it occurs. Their shocks are innovations to the hourly wages of
husbands who keep working, measured as annual earnings over annual hours in two-year changes; a job loss appears as a
fall in hours, and a coincidence of two job losses within a month cannot be seen. High-frequency administrative data
show that this is where most earnings risk lies: \citet{ganong2026earnings} find, with payroll and bank-account records,
that in three quarters of months a worker's pay differs from the month before and in a quarter of months by at least
17\%, that the changes come from hours rather than wages, that about half of a typical worker's hours volatility is caused
by firm-level fluctuations in total hours, that lower-income hourly workers face the most, that the volatility raises
consumption volatility for workers with little liquidity, and that it is masked in annual data. Their analysis stays
within ongoing jobs and studies workers rather than couples; the joint job loss is the household-level counterpart across
jobs, a firm-driven fall of both spouses' hours to zero in the same month. In the panel of
\citet{blundell2016consumption}, a husband is displaced in 4.4\% of couple-years and both spouses in 0.3\% (28
couple-years); of 57 joint displacements among married couples with husbands aged 30--57, 28 reach their estimation file
and 17 contribute a wage moment (Table~\ref{tab:bps_attrition}), most of the rest lost to the exclusion of the
low-income sample, missing lagged observations and the requirement that husbands work in every year. An estimator built
on average covariances of wage, earnings and consumption growth cannot see an event this rare, and their transitory
wage covariance is not carried by displacement years (0.0063 without them). In the interval of a joint displacement the
husband's earnings fall by 0.335 against 0.123 after a single displacement, and the wife's earnings by an arc change
0.44 larger; consumption does not fall in that interval but falls by more in the following one after a joint than after
a single displacement (Table~\ref{tab:bps_event}: $-0.138^{**}$ with the low-income sample, $-0.071$, $p=0.19$, in their
sample). With 17 to 23 joint events this is suggestive; Section~\ref{sec:household} reports the consumption response in
the full PSID for 1999--2021.

Three further features matter for the comparison. The insurance that family labor supply provides in their accounting
is the wife's hours response, a Marshallian cross-elasticity of $-0.75$ identified from wives who work in both years of
the interval, in couples in which the husband keeps working. Applied to a displaced husband's permanent wage loss of
15--20\%, it implies 190--250 more hours a year from wives who already work; on the extensive margin their estimates imply
that 8.4\% of non-working wives enter work per 10\% permanent fall in the husband's wage, or 13--17\% after a displacement
\citep[p.~428]{blundell2016consumption}. The responses to displacement measured here are smaller: 51 hours a year on
average in the PSID and, from 1984, a gross 47 hours offset by 35 hours of joint job losses, and 4.2 (CPS) to 5.4 (SIPP)
percentage points more employment among wives who were not working. The record of a couple ends when the spouse
changes, and the job loss itself changes partners, so a sample of stable couples keeps the households in which the
marriage survived the shock. And the safety net in their schedule is a function of positive family earnings with no
unemployment insurance: the Earned Income Tax Credit pays nothing in the state in which both spouses have lost their
jobs, which lies outside a framework built on positive hours, and the households who depend most on food stamps are
largely excluded by dropping the low-income sample and husbands with a year without earnings, although 71\% of the
insurance of the food-stamp-eligible households who remain comes from the government. The joint job loss puts these
households and this state at the center: it ranks transfers by whether they pay when no one in the household works.

The two sets of results are therefore complements. Family labor supply insures permanent wage shocks among stable,
continuously employed couples, and it fails against the shock at which the second earner's job is lost with the first,
in the month the household's income falls most. A model of family insurance has to deliver both: the model of
Section~\ref{sec:model} contains correlated wage shocks and a progressive tax and transfer schedule as in
\citet{blundell2016consumption}, adds a displacement hazard with a common component, the extensive margin, a search
friction, unemployment insurance and separation, and is required to reproduce their estimates---the spouses' wage-shock
covariances, the split of insurance between family labor supply, government and savings, and how these move when their
restrictions are relaxed---when their sample selection, trims and two-year differencing are applied to simulated
couples, and the joint job loss when they are not.

\paragraph{Job loss, the family and children.} The costs of job loss go beyond earnings, to mortality
\citep{sullivan2009job}, health \citep{black2015losing} and divorce \citep{charles2004job, eliason2012lost}, and they
differ across workers, establishments and markets, with family linkages among the dimensions
\citep{halla2020job, athey2026heterogeneous}. In this literature the unit hit by the shock is one worker, and the
family enters as the spouse who responds. \citet{charles2004job} find in the PSID that a job loss of either the
husband or the wife raises divorce, after layoffs but not after plant closings, and read it as news about the
partner's fitness as a mate. What a job loss of both spouses does to the marriage has not been studied;
Section~\ref{sec:household} confirms their finding for wives over 1984--2021 and shows that a joint loss gives the
largest break-up response. Children whose fathers were displaced by plant closures in Canada earn about 9\% less as adults
\citep{oreopoulos2008intergenerational}; in the PSID the loss is similar with layoffs but close to zero with closures
only \citep{page2009parental}, and it is larger for Black youths' post-secondary education \citep{kalil2011parental}; in
US tax data, fathers' layoffs when children are 12 to 29 sharply reduce family income but have small effects on college
enrollment and early earnings \citep{hilger2016parental}; timing matters \citep{carneiro2022timing}. The PSID evidence in
Section~\ref{sec:children} is consistent with this literature and does not extend it; the paper uses it to bound the
child block of the model.

\paragraph{The safety net as insurance.} Childhood exposure to food stamps and the Earned Income Tax Credit raises
children's long-run outcomes \citep{hoynes2016long, bastian2018long, bailey2024social}. Unemployment insurance is the
largest public buffer against job loss, but the lowest-income job losers are the least insured \citep{east2020how};
more generous unemployment insurance reduces grade repetition among children of job losers \citep{regmi2018examining}
and mitigates the rise in divorce after men's layoffs \citep{lindo2025stable}. These studies vary one program at a time.
The joint job loss ranks programs by whether they pay when no one in the household works.

\paragraph{Structural models of family insurance and mobility.} Dynastic models study how early investment, education
aid and transfers shape intergenerational persistence \citep{lee2019intergenerational, caucutt2020early,
abbott2019education, daruich2018macroeconomic}, and \citet{bolt2024intergenerational} decompose the IGE in estimated
production functions and decision rules. In these models the risk parents face is at most their own earnings risk:
none has a spouse's employment risk or insurance within the household. The models of family insurance that have both
spouses \citep{blundell2016consumption, birinci2019spousal, mankart2017household} do not follow the household to the
children. The model of Section~\ref{sec:model} joins the two with the correlation of the spouses' job losses at its
center.

\begin{table}[!htbp]\centering
\caption{The joint labor-market risk of spouses: what is known and what this paper adds}\label{tab:contribution}
\footnotesize
\begin{tabular}{p{0.38\textwidth}p{0.30\textwidth}p{0.24\textwidth}}
\toprule
Finding & Prior evidence & This paper \\
\midrule
Spouses' earnings covary, mostly through permanent components & \citet{hyslop2001rising}; \citet{shore2010better} & Known \\
\addlinespace
Couples sort on employer and industry; about a tenth share an establishment & \citet{hyatt2019coworking} & Known \\
\addlinespace
Joint displacements by the same plant event exist & \citet{halla2020job} (couples dropped) & Known and named, not measured \\
\addlinespace
Non-participating spouses enter the labor force when the primary earner loses his job, more among young couples; joint
search model with independent separation shocks & \citet{bacher2025joint} & Known; the separation shocks are correlated
here \\
\addlinespace
Correlated unemployment mutes the US added-worker effect & \citet{kawano2025who} & Stated as one of two explanations, not decomposed \\
\addlinespace
Most earnings risk is monthly, from hours not wages, largely firm-driven, highest for low-income hourly workers,
moves consumption, masked in annual data & \citet{ganong2026earnings} & Known for individual workers within jobs; the
joint job loss is its household counterpart across jobs \\
\addlinespace
Family labor supply provides 63\% of the insurance of the husband's permanent wage shocks; spouses' transitory wage
shocks correlated (0.31), permanent ones not & \citet{blundell2016consumption} & Known for wage shocks of continuously employed husbands; reconciled
here with the failure against joint job loss (Section~\ref{sec:literature}) \\
\addlinespace
Monthly within-couple job-loss hazard: tenfold spike in the same month, 0.6--0.8 points in the months around it,
1994--2023, 1.3 million couples, within state $\times$ month $\times$ wife's industry & None & New (Section~\ref{sec:jointloss},
CPS) \\
\addlinespace
Excess displacement with each spouse's own reason for the job ending; future-displacement placebo; her employer closes
when his does; concentration in the same three-digit industry; tetrachoric correlation 0.24, 0.50 in the same sector
& None at this level & New (Section~\ref{sec:jointloss}, PSID) \\
\addlinespace
Timing: her own job loss peaks when the household's loss peaks, while added work arrives later & \citet{stephens2002worker}
documents the build-up of the added-worker response & New: her own loss set next to the added-worker response \\
\addlinespace
Gross versus net added-worker effect & Suggested by \citet{kawano2025who} & New decomposition (Section~\ref{sec:household}) \\
\addlinespace
The excess is equal in both races & None & New \\
\addlinespace
Father's displacement and children's outcomes & \citet{oreopoulos2008intergenerational}; \citet{page2009parental};
\citet{kalil2011parental}; \citet{hilger2016parental} & Replication with a reason-based displacement measure \\
\addlinespace
Divorce after job loss & \citet{charles2004job}; \citet{eliason2012lost} & Known for either spouse; new: a joint loss
of both jobs gives the largest response (the wife's equal effect confirmed over 1984--2021) \\
\addlinespace
Mothers' work around a birth and children's outcomes & \citet{herbst2017are} and \citet{willage2022postpartum} find
harm & Informative nulls (Appendix~\ref{app:other}) \\
\bottomrule
\end{tabular}
\end{table}

The paper does not claim to be the first to show that a father's job loss lowers children's outcomes, that divorce
rises after job loss, that spouses' earnings covary, or that public insurance crowds out spousal labor supply.
